\section{Site Planning and Resilience in Errachidia}
\label{sec:site_planning}

\subsection{Errachidia Site Context}\label{subsec:errachidia-site-context}
Errachidia sits in a semi-arid climate with historic khettara/foggara subterranean water channels.
Our deployment constraints prioritize non-intrusive sensing and authenticity of context:
\begin{itemize}[nosep]
  \item \textbf{Semi-arid conditions:} Large diurnal swings, dust, and seasonal flash-flood risk in wadis inform placement and maintenance windows.
  \item \textbf{Heritage:} Khettara alignments and associated masonry demand reversible mounting (no drilling) and minimal visual/physical impact.
  \item \textbf{Non-intrusive sensing:} Pods focus on RH/T, flood sentinel, and (future) bio-impedance coverage without altering flow or fabric.
\end{itemize}

\subsection{Regulatory Context: Water Domain and Loi 36-15}
\label{subsec:water-law}

Khettara systems are traditional gravity-fed underground water channels that tap shallow
aquifers through a series of vertical shafts and subterranean galleries.
Under Moroccan law, these structures and the water they convey fall within the
\emph{domaine public hydraulique} (public hydraulic domain).

\paragraph{Loi n\textsuperscript{o}\,36-15 relative à l'eau.}
Morocco's water law, promulgated in 2016, establishes the regulatory framework for
water resource management, protection, and monitoring\footnote{\url{https://climate-laws.org/document/law-no-36-15-on-water_f432}}.
Its key principles relevant to TrackOne's deployment include:

\begin{itemize}[nosep]
    \item \textbf{Water as public domain:}
          All water resources, including groundwater and traditional water systems
          (khettaras, foggaras, seguias), are classified as public domain.
          Any intervention in or near these systems---including sensor installation---must
          respect this status and the authority of the relevant basin agency\footnote{\url{https://henry.baw.de/bitstreams/95e5282c-8221-4d8a-b641-5f6b8c89d8ed/download}}.
    \item \textbf{Decentralized basin governance:}
          Management is delegated to regional \emph{Agences de Bassin Hydraulique} (ABH).
    The Errachidia site falls under the jurisdiction of the
    \ac{ABH} Guir-Ziz-Rhéris.
    Deployment authorization, if required for sensor placement in or near the
    hydraulic domain, would be sought from this agency.
    \item \textbf{Monitoring and information systems:}
          The law mandates the establishment of water-related information systems,
          including monitoring of water quality and quantity\footnote{\url{https://water.fanack.com/morocco/water-management-in-morocco/}}.
          TrackOne's verifiable telemetry pipeline---producing tamper-evident,
          timestamped environmental records---directly supports this objective.
    \item \textbf{Protection and conservation:}
          The law strengthens mechanisms for protecting water resources from pollution
          and degradation.
          Continuous RH/T monitoring and flood-sentinel alerting in khettara galleries
          contribute to early detection of conditions that may indicate structural
          deterioration, infiltration changes, or flood risk.
    \item \textbf{User-pays / polluter-pays:}
          The law establishes financial instruments based on usage and pollution fees.
          Verifiable, anchored sensor data could serve as an evidentiary basis for
          regulatory reporting under this framework.
\end{itemize}

\paragraph{Alignment with TrackOne's design principles.}
Several design choices in TrackOne were made with this regulatory context in mind:

\begin{table}[H]
\centering
\small
\caption{Alignment between Loi 36-15 requirements and TrackOne design properties.}
\label{tab:law-36-15-alignment}
\begin{tabular}{@{}p{4.2cm}p{5.0cm}p{4.2cm}@{}}
\toprule
\textbf{Loi 36-15 requirement}     & \textbf{TrackOne mechanism}                           & \textbf{Reference} \\
\midrule
Water monitoring systems            & Continuous RH/T and flood-sentinel telemetry           & FR-2, \ac{ADR}-027                                \\
Data integrity for regulatory use   & \ac{AEAD} + Merkle + \ac{OTS} anchoring (tamper-evident records) & \ac{ADR}-003, \ac{ADR}-024 \\
Non-intrusive intervention          & Reversible pod mounting; no drilling or flow alteration & \S\ref{sec:site_planning} \\
Public domain respect               & Deployment within approved polygons; \ac{GIS}-validated placement & \S\ref{subsec:traceability} \\
Decentralized governance support    & Independent verifiability; data readable without trust in operator & Architecture \S\ref{subsec:arch-responsibilities} \\
\bottomrule
\end{tabular}
\end{table}

\paragraph{Practical implications for deployment.}
\begin{itemize}[nosep]
    \item \textbf{Authorization:}
    Prior to production deployment, the project must determine whether \ac{ABH}
    Guir-Ziz-Rhéris requires a formal declaration, light-licensing,
    or authorization for sensor placement within the hydraulic domain.
    The non-intrusive, reversible nature of pod installation is expected to
    simplify this process.
          % TODO [LAW-1]: Contact ABH Guir-Ziz-Rhéris to confirm the authorization
          %   procedure for non-intrusive sensor deployment in khettara galleries.
          %   Document the procedure and any conditions imposed.
    \item \textbf{Data sovereignty:}
          Loi36--15implies that monitoring data related to the public hydraulic domain
          may be subject to reporting obligations or requests from the basin agency.
          TrackOne's append-only ledger and independent verifiability ensure that any
          data shared with regulators is demonstrably unaltered.
    \item \textbf{Heritage considerations:}
          While Loi36--15governs the water domain, khettaras may also be subject to
          cultural heritage protections under separate legislation.
          Coordination with both water and heritage authorities is prudent.
          % TODO [LAW-2]: Check whether the target khettara segments are listed under
          %   any cultural heritage protection (e.g., Loi 22-80 on conservation of
          %   historic monuments). If so, note any additional constraints.
\end{itemize}

\subsection{GIS Layers and Placement Methodology}
\label{subsec:gis-placement}
Planning relies on OpenStreetMap/GEOFABRIK layers and \ac{DEM} overlays:
\begin{itemize}[nosep]
  \item \textbf{OSM layers (Geofabrik extract):} water, waterways, landuse, natural, roads, places.
  \item \textbf{Derived overlays:} Khettara alignments/channel networks; elevation (\ac{DEM} from \ac{SRTM} or similar) for flood basin inference.
  \item \textbf{Method:} Overlay alignments with land use and access; score candidate sites by coverage, accessibility, and redundancy; select from pre-approved polygons and route maintenance accordingly.
\end{itemize}
Coordinates and polygons are versioned alongside deployment manifests.

\subsection{Frequency Selection and Regulatory Constraints}
\label{subsec:frequency-regulatory}

The choice of 868\,MHz \ac{ISM} band for TrackOne's \ac{LoRa} uplink is driven by a combination of
propagation physics, hardware availability, and Moroccan spectrum regulation.

\paragraph{Propagation rationale.}
Sub-GHz frequencies offer significantly better penetration through earth, masonry, and
moisture-laden materials than 2.4\,GHz alternatives (Wi-Fi, BLE, Zigbee).
For subterranean khettara galleries with mortared stone walls and variable soil cover,
868\,MHz provides a favorable tradeoff between antenna size, link budget, and
material attenuation.
\ac{LoRa}'s chirp spread spectrum (CSS) modulation adds processing gain
(up to $\sim$20\,dB at \ac{SF}12) that further extends usable range in high-attenuation
environments.

\paragraph{Moroccan regulatory framework.}
Radio frequency spectrum in Morocco is administered by the
\emph{Agence Nationale de Réglementation des Télécommunications} (ANRT)
under the authority of \emph{Loi n\textsuperscript{o}\,36-15} (2015), which establishes
the general framework for telecommunications regulation, including spectrum allocation,
licensing, and enforcement.
% TODO [REG-1]: Cite the official Gazette publication of Loi 36-15:
%   Dahir n° 1-16-25 du 19 rejeb 1437 (27 avril 2016) portant promulgation de la
%   loi n° 36-15 relative aux télécommunications. Bulletin Officiel n° 6470.

\paragraph{ANRT channel plan for 868\,MHz.}
\ac{ANRT}'s national frequency allocation table designates the 868--870\,MHz band for
short-range devices (\ac{SRD}s) under conditions broadly aligned with European
harmonisation (CEPT/ECC Recommendation70--03 ETSI EN\,300\,220):

\begin{table}[H]
\centering
\small
\caption{ANRT 868\,MHz sub-band constraints relevant to LoRa SRDs.}
\label{tab:anrt-868}
\begin{tabular}{@{}llll@{}}
\toprule
\textbf{Sub-band}       & \textbf{Frequency}       & \textbf{Max ERP}  & \textbf{Duty cycle} \\
\midrule
g1 (general \ac{SRD})        & 868.0--868.6\,MHz        & 25\,mW (14\,dBm)  & $\le$\,1\,\%        \\
g2 (general \ac{SRD})        & 868.7--869.2\,MHz        & 25\,mW (14\,dBm)  & $\le$\,0.1\,\%      \\
g3 (high duty cycle)    & 869.4--869.65\,MHz       & 500\,mW (27\,dBm) & $\le$\,10\,\%       \\
\bottomrule
\end{tabular}
\end{table}

% TODO [REG-2]: Verify these exact sub-band limits against the current ANRT
%   "Tableau National d'Attribution des Fréquences" (TNAF) and any SRD-specific
%   Decision. ANRT periodically updates these; confirm the version in force as of
%   2026. If Morocco deviates from CEPT/ECC 70-03 on any sub-band, note the
%   difference explicitly.
% TODO [REG-3]: If ANRT requires a declaration or light-licensing for SRD use
%   (even if no individual license is needed), document the procedure here.

\paragraph{Implications for TrackOne duty cycling.}
The 1\,\% duty-cycle limit in sub-band g1 is the binding constraint for TrackOne's
uplink cadence.
At \ac{SF}9/BW\,125\,kHz (a likely operating point), a single 50-byte \ac{LoRa} frame requires
$\sim$200\,ms of air time, permitting approximately 180 transmissions per hour within
the 1\,\% budget.
TrackOne's default cadence of 6 frames/hour (one every 10 minutes) consumes
$<$\,0.04\,\% of duty cycle---well within the regulatory envelope and leaving substantial
headroom for retransmissions, adaptive bursts (\ac{ADR}-025), and future OTA command channels
(\ac{ADR}-026).

\begin{table}[H]
\centering
\small
\caption{Duty-cycle budget at TrackOne's default uplink cadence (sub-band g1).}
\label{tab:duty-cycle-budget}
\begin{tabular}{@{}lrl@{}}
\toprule
\textbf{Parameter}            & \textbf{Value}       & \textbf{Source} \\
\midrule
Regulatory duty-cycle limit   & 1\,\%          & \ac{ANRT} / CEPT g1                   \\
Frame size (\ac{AEAD} payload)     & $\le$\,50\,B         & \ac{ADR}-002                          \\
Air time per frame (\ac{SF}9/125)  & $\sim$200\,ms        & SX1276 datasheet \\
Default cadence               & 6 frames/hr          & \ac{ADR}-025                          \\
Consumed duty cycle           & $\sim$0.03\,\%       & $6 \times 0.2\text{s} / 3600\text{s}$ \\
Remaining headroom            & $>$\,96\,\%          & Available for retx / OTA / burst \\
\bottomrule
\end{tabular}
\end{table}

\paragraph{Why not 433\,MHz or 915\,MHz?}
\begin{itemize}[nosep]
    \item \textbf{433\,MHz:} Better penetration but significantly larger antenna
          ($\lambda/4 \approx 17$\,cm vs.\ 8.6\,cm at 868\,MHz), complicating
    non-intrusive pod enclosure design.
    \ac{ANRT} allocation for 433\,MHz \ac{SRD}s
    is narrower with stricter power limits.
    \item \textbf{915\,MHz:} Not available in Morocco (nor in CEPT Region~1).
          The 902--928\, MHz band is allocated to ITU Region~2 (Americas);
    using it in Morocco would violate \ac{ANRT} regulations.
    \item \textbf{2.4\,GHz (BLE/Wi-Fi/Zigbee):} Unsuitable propagation through
    earth and masonry; link budget inadequate for subterranean deployment
    without dense relay infrastructure.
\end{itemize}

\paragraph{Why raw LoRa, not LoRaWAN.}
TrackOne uses the SX1276 \ac{LoRa} physical layer (CSS modulation) directly,
without the \ac{LoRa}WAN MAC and network stack.
This is a deliberate choice
driven by four constraints:

\begin{enumerate}[nosep]
        \item \textbf{Cryptographic sovereignty:}
        \ac{LoRa}WAN hardcodes AES-128 for session encryption and integrity
        with no algorithm negotiation.
        TrackOne requires XChaCha20-Poly1305 (256-bit keys, 192-bit nonces)
        and algorithm agility for post-quantum readiness
        (see~\S\ref{subsec:pq-risk-growth}).
        The \ac{LoRa} Alliance roadmap places crypto-agility as a future
        development item (2024--2027 phase), but no specification has been
        published as of early 2026.
        \item \textbf{OTA firmware update model:}
        \ac{LoRa}WAN's FUOTA (Firmware Update Over The Air) requires Class~C
        continuous-receive mode during update sessions, multicast group
        management via the network server (TS005), fragmentation (TS004),
        and clock synchronization (TS003).
        This imposes significant power cost on duty-cycled, energy-harvesting
        pods and creates a dependency on a \ac{LoRa}WAN Network Server that
        TrackOne's architecture deliberately avoids.
        TrackOne instead uses a staged bootloader with Ed25519-signed firmware
        images delivered over the existing authenticated uplink/downlink
        channel (\ac{ADR}-026), requiring no additional protocol layers or
        infrastructure.
        \item \textbf{Evidence integrity independence:}
        \ac{LoRa}WAN delegates trust to the Network Server for key management,
        session state, and frame-counter enforcement.
        TrackOne's threat model requires that evidence integrity
        (\ac{AEAD}~+~Merkle~+~\ac{OTS} anchoring) be verifiable \emph{without trusting
        any network infrastructure} (see Architecture,
        \S\ref{subsec:arch-responsibilities}).
        Introducing a \ac{LoRa}WAN Network Server would add a trusted third party
        that the system is explicitly designed to avoid.
        \item \textbf{Deployment simplicity:}
        A raw \ac{LoRa} link between pod and gateway requires no network server,
        no join server, no application server, and no \ac{LoRa}WAN certification.
        For a research deployment of fewer than 20 pods at a single site,
        the operational overhead of a \ac{LoRa}WAN stack is not justified.
\end{enumerate}

\paragraph{Future LoRaWAN alignment.}
This choice is \emph{not} a permanent rejection of \ac{LoRa}WAN.
If a future \ac{LoRa}WAN revision delivers crypto-agility with an algorithm menu
that includes 256-bit \ac{AEAD} and negotiable key exchange, and if FUOTA evolves
to support duty-cycled Class~A devices without continuous-receive windows,
then adopting the \ac{LoRa}WAN MAC would be reconsidered.
TrackOne's layered architecture (crypto above transport,
\S\ref{subsec:layering-wireguard}) is designed so that the transport layer
can be replaced without affecting the evidence pipeline.

\subsection{RF Characterization Status}
\label{subsec:rf-characterization}

A preliminary bench test was conducted with two LILYGO TLORA V2.1 nodes
(SX1276, 868\,MHz, Meshtastic firmware) to validate radio hardware functionality.
\ac{LoRa} TX/RX was confirmed at close range (RSSI\,$= -28$\,dBm, SNR\,$= 6.25$\,dB),
but operational errors (modem preset mismatch, app-to-radio path confusion)
prevented collection of subterranean propagation data during the field visit.

A follow-up measurement campaign with corrected procedures (serial debug logging,
appropriate spreading factor, separated configuration, and measurement phases)
is planned to produce quantitative link-margin data at multiple gallery positions.
Results will inform spreading-factor selection and the adaptive uplink cadence
policy defined in \ac{ADR}-025.
% TODO [RF-FOLLOWUP]: After follow-up campaign, expand this subsection into a
%   proper link-budget characterization with RSSI/SNR vs. position table and
%   SF-range tradeoff analysis.

% ---------------------------------------------------------------------------
% Figure A2 – Micro-site DEM cross-section
\begin{figure}[H]
  \centering
  %\includegraphics[width=0.9\textwidth]{figs/microsite_dem_cross_section.png}
  \fbox{\parbox[c][3.0cm][c]{0.9\linewidth}{Micro-site DEM cross-section placeholder (elevation vs distance with pod depth and gallery annotations)}}
  \caption{Elevation cross-section across a khettara segment showing pod depths, gallery location, and above-ground gateway placement.}
  \label{fig:dem-cross-section}
\end{figure}

\subsection[Resilience and Quality Attributes (ISO 25010)]{Resilience and Quality Attributes (ISO~25010)}\label{subsec:resilience-and-quality-attributes-(iso25010)}
We tie siting choices to relevant quality attributes:
\begin{itemize}[nosep]
  \item \textbf{Reliability:} At least two pods per critical segment; redundant coverage across bends and RF shadows.
  \item \textbf{Security:} Pods placed in protected locations with clear gateway and calendar placement; verifiability maintained via stationary calendar when public pools are unstable.
  \item \textbf{Maintainability:} Avoid inaccessible shafts; \ac{GIS} planning favors reachable sites and efficient maintenance routes.
\end{itemize}
These attributes link to requirements in Section~\ref{sec:system-requirements} (FR-2, FR-6, RR-5; NFR-1/NFR-2) and to the safety net in Section~\ref{sec:security}.

\subsection{Traceability from Plan to Deployment}
\label{subsec:traceability}
Validation ties pod IDs to \ac{GIS} coordinates:
\begin{itemize}[nosep]
  \item Pod IDs are bound to coordinates within approved polygons; a simple \ac{GIS} script checks that no pods are orphaned or outside legal regions.
  \item Pre-deployment checks ensure coordinates fall within the designated polygons and that maintenance paths exist.
  \item This validation is part of the release pipeline and is traceable back to the requirements matrix (Section~\ref{subsec:requirements-traceability}).
\end{itemize}
